\documentclass[../../../include/open-logic-section]{subfiles}

\begin{document}

\olfileid{sth}{replacement}{extrinsic}
\olsection[ملاحظات بیرونی]{ملاحظات بیرونی دربارهٔ جایگزینی}

با بررسی کوششی \emph{بیرونی} برای توجیه جایگزینی آغاز می‌کنیم. بولوس
پیشنهاد زیر را مطرح می‌کند:
\begin{quote}
  [\ldots] دلیل پذیرش اصول موضوع جایگزینی بسیار ساده است: آن‌ها پیامدهای
  مطلوب فراوانی دارند و (ظاهراً) هیچ پیامد نامطلوبی ندارند. افزون بر
  قضایای مربوط به تصور تکراری، این پیامدها نظریه‌ای رضایت‌بخش، هرچند نه
  آرمانی، دربارهٔ اعداد نامتناهی و نتیجه‌ای بسیار مطلوب را دربر می‌گیرند
  که تعریف‌های استقرایی روی روابط خوش‌بنیاد را توجیه می‌کند.
  \citep[229]{Boolos1971}
\end{quote}
مغز اندیشهٔ بولوس آن است که جایگزینی را با ثمراتش توجیه کنیم. ثمرات
خاصی که او ذکر می‌کند همان چیزهایی‌اند که در چند فصل گذشته بررسی
کرده‌ایم. جایگزینی به ما اجازه داد اثبات کنیم اعداد ترتیبی فون نویمان
جانشین‌هایی عالی برای مفهوم نوعِ خوش‌ترتیب‌اند (این همان «نظریهٔ
رضایت‌بخش، هرچند نه آرمانی، دربارهٔ اعداد نامتناهی» است). جایگزینی همچنین
به ما اجازه داد~$V_\alpha$ها را تعریف کنیم، مفهوم رتبه را بنا نهیم و
استقرای~$\in$ را اثبات کنیم (این‌ها «قضایای مربوط به تصور تکراری» را
تشکیل می‌دهند). سرانجام، جایگزینی به ما اجازه می‌دهد قضیهٔ بازگشت
ترامتناهی را اثبات کنیم (این همان «تعریف‌های استقرایی روی روابط
خوش‌بنیاد» است).

این‌ها واقعاً پیامدهایی مطلوب‌اند. اما آیا همین پیامدهای مطلوب برای
\emph{توجیه} جایگزینی کافی‌اند؟ \emph{نه}. یا دست‌کم نه به‌طور
سرراست.

در اینجا مشکلی ساده وجود دارد. هرچند چند پیامد مطلوب جایگزینی را بیان
کرده‌ایم، می‌توانستیم بسیاری از آن‌ها را از راه‌های دیگری نیز به دست
آوریم. این امر آن‌قدر که باید شناخته‌شده نیست؛ پس شایسته است اندکی درنگ
کنیم و وضع را توضیح دهیم.

نظریهٔ ساده‌ای دربارهٔ مجموعه‌ها به نام نظریهٔ ترازها، یا به‌اختصار
$\LT$، وجود دارد.\footnote{نخستین صورت‌های~$\LT$ را
\citet{Montague1965} و~\citet{Scott1974} عرضه کردند؛
\citet{Potter2004} آن را ساده‌تر کرد و شرحی در اندازهٔ یک کتاب از آن
به دست داد؛ و~\citet{ButtonLT1} به‌تازگی~$\LT$ را باز هم ساده‌تر کرده
است.} اصول موضوع~$\LT$ فقط گسترش‌مندی، جداسازی و این ادعا هستند که هر
مجموعه زیرمجموعهٔ یک \emph{تراز} است؛ «تراز» نیز زیرکانه چنان تعریف
می‌شود که ترازها مانند دوستان ما، $V_\alpha$ها، رفتار کنند. پس~$\ZF$،
$\LT$ را اثبات می‌کند؛ اما~$\LT$ \emph{بسیار} ضعیف‌تر از~$\ZF$ است.
در واقع، $\LT$ زوج‌ها، مجموعه‌های توانی، نامتناهی یا جایگزینی را به دست
نمی‌دهد. بگذارید~$\Zr$ حاصل افزودن نامتناهی و مجموعه‌های توانی به~$\LT$
باشد؛ این نظریه زوج‌ها را نیز به دست می‌دهد، پس~$\Zr$ دست‌کم به‌اندازهٔ
$\Z$ قوی است. اما در واقع~$\Zr$ اکیداً قوی‌تر از~$\Z$ است، زیرا این
ادعا را می‌افزاید که هر مجموعه رتبه‌ای دارد (ازاین‌رو پیشنهاد می‌کنم آن
را~$\Zr$ بنامیم). در واقع، $\Zr$ این‌ها را به دست می‌دهد: نظریه‌ای
کاملاً رضایت‌بخش دربارهٔ اعداد ترتیبی؛ نتایجی که سلسله‌مراتب را به مراحل
خوش‌ترتیب لایه‌بندی می‌کنند؛ اثبات استقرای~$\in$؛ و \emph{صورتی} از
بازگشت ترامتناهی.

خلاصه آنکه، هرچند بولوس از این امر آگاه نبود، همهٔ پیامدهای مطلوبی که
ذکر می‌کند می‌توانستند \emph{بدون} جایگزینی به دست آیند؛ کافی بود به‌جای
$\Z$ از~$\Zr$ استفاده کند.

(با توجه به همهٔ این مطالب، چرا به‌جای آموزش~$\LT$ و~$\Zr$، مسیر متعارف
آموزش~$\ZF$ را در پیش گرفتیم؟ دو دلیل دارد. نخست، صرفاً به دلایل تاریخی،
آغازکردن با~$\LT$ بسیار نامتعارف است؛ می‌خواستیم شما را برای خواندن
بحث‌های معیارتر دربارهٔ نظریهٔ مجموعه‌ها آماده کنیم. دوم، هنگامی که برای
درک~$\LT$ و~$\Zr$ آماده شدید، می‌توانید به‌سادگی
\citealt{Potter2004} و~\citealt{ButtonLT1} را بخوانید.)

البته چون~$\Zr$ اکیداً ضعیف‌تر از~$\ZF$ است، نتایجی وجود دارند که~$\ZF$
آن‌ها را اثبات می‌کند اما~$\Zr$ درباره‌شان ساکت می‌ماند. بنابراین ممکن
است کسی بکوشد جایگزینی را بر مبنایی بیرونی با اشاره به یکی از این نتایج
توجیه کند. اما وقتی کاربرد~$\Zr$ را بدانید، یافتن نمونه‌های فراوانی از
چیزهایی که (الف) جایگزینی درباره‌شان تعیین تکلیف می‌کند ولی راه دیگری
برای تعیین تکلیفشان نیست، و (ب) به‌طور شهودی صادق‌اند، بسیار دشوار است.
(برای مطالب بیشتر بنگرید به~\citealt[\S13.2]{Potter2004}.)

حاصل سخن چنین است. برای عرضهٔ توجیهی بیرونی و قانع‌کننده برای جایگزینی،
باید نتیجه‌ای یافت که \emph{بدون} جایگزینی قابل حصول نباشد. و این کار
آسانی نیست.

اکنون مشکل دیگری را بررسی کنیم که در برابر هر کوشش برای ارائهٔ توجیهی
صرفاً بیرونی برای جایگزینی پدید می‌آید. (شاید این مشکل از مشکل نخست
بنیادی‌تر باشد.) بولوس فقط اشاره نمی‌کند که جایگزینی پیامدهای مطلوب
فراوانی دارد؛ او همچنین می‌گوید جایگزینی «(ظاهراً) هیچ پیامد نامطلوبی»
ندارد. اما این قید معترضه، یعنی «ظاهراً»، بی‌گمان کاملاً اساسی است.

به یاد آورید چگونه به اینجا رسیدیم: اصل موضوع فهم خام به ناسازگاری
انجامید، و ما با پذیرفتن تصور انباشتی-تکراریِ مجموعه به آن پاسخ دادیم.
این تصور با روایتی همراه است که امیدواریم ما را از سازگاری‌اش مطمئن
سازد. اما اگر نتوانیم جایگزینی را از درون آن روایت توجیه کنیم، هنوز هیچ
دلیلی برای باور به سازگاری~$\ZF$ نداریم. یا دقیق‌تر بگوییم: جز این واقعیت
(شاید صرفاً تصادفی) که \emph{هنوز} هیچ‌کس تناقضی کشف نکرده است، دلیلی
برای باور به سازگاری~$\ZF$ نداریم. ملاحظهٔ بولوس دقیقاً به این معنا
تقلیل می‌یابد: «(ظاهراً)~$\ZF$ سازگار است». باید اطمینانی بیش از این
دربارهٔ سازگاری طلب کنیم.

این مسئله هر کوشش \emph{صرفاً} بیرونی برای توجیه جایگزینی را متأثر
می‌سازد؛ یعنی هر توجیهی که تنها بر حسب پیامدهای شناخته‌شدهٔ~$\ZF$ بیان
شود. بنابراین در پی توجیهی \emph{درونی} برای جایگزینی خواهیم بود؛ یعنی
توجیهی که نشان دهد روایتی که دربارهٔ مجموعه‌ها گفتیم، به‌نحوی از پیش ما
را به جایگزینی ملتزم می‌کند.

\end{document}
